\documentclass[tikz,border=8pt]{standalone}
\usepackage{ctex}
\usepackage{amssymb}
\usepackage{pgfplots}
\pgfplotsset{compat=1.18}
\usetikzlibrary{calc,positioning,arrows.meta}

\begin{document}
\begin{tikzpicture}

%% ===== 左图：凸集（椭圆） =====
\begin{scope}[xshift=0cm]

  % 填充椭圆
  \filldraw[gray!20, draw=black, thick] (0,0) ellipse (2.5 and 1.8);

  % 取两个内部点
  \coordinate (A) at (-1.2, 0.5);
  \coordinate (B) at (1.4, -0.6);

  % 连线（线段在集合内）
  \draw[red!80!black, thick] (A) -- (B);

  % 点
  \fill[blue!70!black] (A) circle (3pt) node[above left, font=\small] {$\mathbf{x}$};
  \fill[blue!70!black] (B) circle (3pt) node[below right, font=\small] {$\mathbf{y}$};

  % 中点标注
  \coordinate (M) at ($(A)!0.5!(B)$);
  \fill[red!60!black] (M) circle (2.5pt);
  \node[above, font=\scriptsize, red!60!black] at (M) {$\theta\mathbf{x}+(1{-}\theta)\mathbf{y}$};

  % 右箭头注释（在集合内）
  \node[draw=gray!50, fill=white, thin, rounded corners=2pt,
        font=\scriptsize, inner sep=3pt] at (0, -2.5)
    {线段 $\subset C$：凸集};

  % 标题
  \node[font=\small\bfseries] at (0, 2.4) {凸集};

  % 验证公式
  \node[font=\scriptsize] at (0, -3.3)
    {$\forall\,\mathbf{x},\mathbf{y}\in C,\,\theta\in[0,1]:\,\theta\mathbf{x}+(1-\theta)\mathbf{y}\in C$};

\end{scope}

%% ===== 右图：非凸集（月牙形） =====
\begin{scope}[xshift=7.5cm]

  % 月牙形 = 大圆 minus 小圆（用 clip + 填充）
  % 大椭圆
  \begin{scope}
    \clip (-2.8,-2.0) rectangle (2.8,2.0);
    \filldraw[gray!20, draw=black, thick] (0,0) ellipse (2.5 and 1.8);
    % 挖去右侧圆（制造凹陷）
    \filldraw[white, draw=white] (1.0, 0) circle (1.5);
  \end{scope}
  % 重新画边界
  \draw[black, thick] (0,0) ellipse (2.5 and 1.8);
  \draw[black, thick] (1.0, 0) circle (1.5);

  % 取两个集合内的点（月牙形两端）
  \coordinate (C) at (-1.8, 0.8);
  \coordinate (D) at (-0.2, -1.4);

  % 线段穿出集合
  \draw[red!80!black, thick, dashed] (C) -- (D);

  % 线段中某点在集合外的标注
  \coordinate (Out) at (-1.0, -0.3);
  \fill[orange!80!black] (Out) circle (2.5pt);
  \node[right, font=\scriptsize, orange!70!black] at (Out) {在集合外！};

  % 两端点
  \fill[blue!70!black] (C) circle (3pt) node[above left, font=\small] {$\mathbf{x}$};
  \fill[blue!70!black] (D) circle (3pt) node[below right, font=\small] {$\mathbf{y}$};

  % 注释框
  \node[draw=gray!50, fill=white, thin, rounded corners=2pt,
        font=\scriptsize, inner sep=3pt] at (0, -2.5)
    {线段 $\not\subset C$：非凸集};

  % 标题
  \node[font=\small\bfseries] at (0, 2.4) {非凸集（月牙形）};

  % 非凸原因
  \node[font=\scriptsize] at (0, -3.3)
    {存在 $\mathbf{x},\mathbf{y}\in C$，使得 $\theta\mathbf{x}+(1-\theta)\mathbf{y}\notin C$};

\end{scope}

%% 整体标题
\node[font=\large\bfseries] at (7.5, 3.1) {凸集（左）vs 非凸集（右）};

\end{tikzpicture}
\end{document}
